\section{Test Budget Allocation}
\label{sec:alloc}



A challenge faced by the approach of \textbf{causal stitching} 
is the vast number of combinations of faults and workloads. 
For example, HDFS has about 3,000 source code locations 
where exceptions can be thrown 
and 2,000 loops where delays can be injected 
and thousands of integration tests. 
Combining injected faults 
with different workloads can result in
tens of millions of possible fault injection scenarios. 
In this section, 
we explain \tool's \textbf{three-phase allocation (3PA) protocol} of test budget, 
which maximizes the discovery of causal relationships 
within a fixed test budget. 
\footnote{A formal definition of the 3PA protocol is in \Cref{app:3pa}.}

\subsection{Principles of Test Allocation}

The 3PA protocol 
leverages two principles  
to prioritize faults based on 
the uniqueness and diversity of their causal impact: 
1) injecting \textit{causally equivalent} faults to diverse workloads, 
and 2) extending \textit{conditional} causal relationships. 


\textbf{Causally Equivalent Fault}. 
Since 
different faults (e.g., $f_1$ and $f_2$) 
may trigger similar interferences (i.e., additional faults), 
we consider such faults ($f_1$ and $f_2$) \textit{causally equivalent} with respect to their consequences. 
For example, in a try-catch block, 
multiple exceptions in the try block 
might be handled by the same catch block, 
in which triggers an additional fault. 
In such cases, 
it is redundant to inject \textit{causally equivalent} faults 
into the same workload. 
Instead, 
\tool prioritizes to inject them in different workloads 
in order to maximize the discovery of diverse fault propagations. 

\textbf{Conditional Causal Relationship}. 
Because
\vcs 
are often formed by 
fault propagations activated by specific conditions, 
\tool 
allocates test budget 
to increase the probability 
to extend \textit{conditional} causal relationships between faults. 
\tool detects 
such \textit{conditional} causal relationships 
by checking if a fault causes different faults in different workloads. 
For example, 
if $f_1$ always causes $f_2$ in different workloads, $f_1 \rightarrow f_2$ is considered \textit{unconditional}, 
while if $f_1$ causes $f_2$ in one workload but $f_3$ in another, 
then $f_1 \rightarrow f_2$ and $f_1 \rightarrow f_3$ are considered 
\textit{conditional}. 
\tool prioritizes to allocate test budgets 
to trigger the resulting faults ($f_2$ and $f_3$) 
of such conditional causal relationships. 


\subsection{Three-Phase Allocation Protocol}
\label{sec:alloc:3pa}

Following the above principles, 
\tool allocates the test budget in three phases. 
In phase one (\textbf{causally equivalent fault detection}), 
\tool conducts a preliminary exploration 
by running each fault injection against one test workload, 
and clusters faults causing similar interferences 
in order to detect \textit{causally equivalent} faults. 
In phase two (\textbf{causality exploration}), 
\tool picks a fault from each set of \textit{causally equivalent} faults 
and injects it into diverse workloads 
to explore its causal relationships with other faults. 
In phase three (\textbf{conditional-causality-guided causality extension}),  
\tool uses the result from phase two to detect 
\textit{conditional} causal relationships, 
and prioritizes to inject faults that 
can extend such \textit{conditional} causal relationships. 
The total test budget is calculated 
using the number of fault locations in each system, 
currently as $4 \times |\mathbb{F}|$, 
where $\mathbb{F}$ is the set of faults.
Phase one makes up the first 25\% of the budget. Phase two  gets 50\% of the budget, while phase three uses the remaining 25\%.



\textbf{Phase One: Causally Equivalent Fault Detection}. 
In this phase, 
for each fault $f_i$, 
\tool picks the workload $t_{i_1}$ that reaches $f_i$'s program location 
and has the highest code coverage. 
The rationale is that 
such a test is more likely to 
reveal the most diverse interferences of $f_i$. 
Different faults may be injected into different workloads since a single test may not cover all fault locations. 
\tool then performs the fault causality analysis 
to obtain all additional faults triggered. 
We use $I(f_i, t_j)$ to represent a list of additional faults triggered by injecting $f_i$ to $t_j$. 


Injections that do not affect system execution 
will be deprioritized automatically.
These non-impactful injections will be clustered together, and assigned the lowest allocation weight in the remaining phases. 
This complements our conservative static filtering of the injection points (\Cref{sec:fia:faults-to-inject}).  


\tool uses an IDF-based
(inverse document frequency \cite{schutze2008introduction}, more details in \Cref{app:3pa:idf}) 
clustering algorithm 
to detect \textit{causally equivalent} faults.  
IDF is commonly used in text mining tasks 
for applying weights to the words~\cite{beel2016paper}. 
It excels at reducing noises induced by common words such as ``a'' and ``the''. 
\tool uses IDF to assign a weight to each fault $f$ in the fault corpus $\mathbb F$. 
The intuition is that, similar to common words in text mining, 
if a fault is frequently triggered by various other faults (e.g., inside a utility function), it is less useful for determining the similarity of interferences $I(f_i, t_j)$ from different injections. 
\tool uses such weights to determine the similarity of the interferences caused by different injections. 

Specifically, 
\tool vectorizes each 
$I(f_i, t_{i_1})$ as 
$V(f_i, t_{i_1})=\mathrm{IDF_v}(I(f_i, t_{i_1}), \mathbb{F})$, 
resulting in a real vector of length $n$, with each element ranging from 0 to 1 (i.e., $V \in [0,1]^n$).
\tool then performs 
a hierarchical clustering~\cite{kaufman2009finding} of the faults 
in $\mathbb F$ using the cosine distance~\cite{schutze2008introduction} between all vectorized interferences $V(f_i, t_{i_1})$. 
We select hierarchical clustering due to their explainable nature, 
as required in biology~\cite{krause2005large} and medical~\cite{zhang2017hierarchical} researches. 
We use cosine distance instead of euclidean distance due to their insensitive to vector length (i.e., the number of additional faults triggered in this case). The cosine distance here ranges from 0 to 1 because all IDF vectors are positive.

By the end of phase one, each fault $f_i$ is clustered in to a group $G_j$ with other faults having similar interferences on the system once injected.
This phase 
makes up 25\% of the test budget, 
providing an initial understanding of each fault's interference.



\textbf{Phase Two: Causality Exploration.} 
In phase two, 
\tool distributes test budgets among 
fault clusters $G_i$ in a round-robin manner 
to perform injection and explore their inteferences.  
This approach allocates equal test budgets to fault clusters, 
not individual faults, 
avoiding redundant injections that yield similar interferences. 

In both phase two and three, 
each time a cluster receives a test quota, 
\tool \textbf{randomly} selects one fault from that cluster 
and injects it into a new workload. 
This is to compensate for potential inaccuracies in phase one 
due to limited test budget. 
For example, 
injecting $f_1$ and $f_2$ both causes $f_3$ in phase one, 
but $f_1$ can cause $f_4$ in 
another workload 
which is unfortunately not used in phase one. 
Using random selection, every fault within a cluster have a chance be injected into a new workload.

After completing all the injection runs, \tool performs fault causality analysis again to identify additional faults $I(f_i, t_j)$. 
A second IDF vectorizer is trained with the data from both phases to convert $I(f_i, t_j)$ to $V(f_i, t_j)$, 
which is used in prioritization in the next phase. 

\tool prioritizes to 
inject faults that can cause conditional interferences 
to more workloads, 
in order to extend conditional causal relationships. 
It measures the 
\textit{diversity} 
of the causal consequences (interferences) 
of faults in a cluster 
using 
an intra-cluster 
interference similarity score ($\mathrm{SimScore}(G_i)$). 
\tool prioritizes to 
inject faults from clusters with a lower score, 
which increases the chance for the injection to trigger conditional interferences. 
The score is the average pairwise cosine distance between all vectorized interference results $V(f_i, t_j)$ obtained from the injection experiments conducted in cluster $G_i$ during phase one and two. 
$\mathrm{SimScore}(G_i)$ ranges from 0 to 1. 
A value of 1 indicates that all faults in cluster $G_i$ trigger the same set of additional faults among all injection runs.
The score is then used in phase three. 



\textbf{Phase Three: Conditional-Causality-Guided Causality Extension.}
Using the $\mathrm{SimScore}(G_i)$ calculated in phase two, 
\tool performs a weighted random allocation favoring clusters 
with more conditional causal consequences. 
Lower similarity indicates 
more conditional causal consequences 
and thus a higher allocation weight.
The allocation weight for cluster $G_i$ is defined as 
$\max{(\epsilon, 1-\mathrm{SimScore}(G_i))}$.
Each cluster has a minimum weight $\epsilon$ of 0.01, ensuring every cluster receives some budget, even with perfectly matched intra-group interference results.


\textbf{Test Budget Transfer.} 
In phase two and three, 
test quotas are \textit{transferable} between clusters. 
In phase two, 
if a cluster exhausts its (fault, test) combinations before using up its quota, the remaining quota is randomly transferred to a larger cluster for more thorough exploration.
In phase three, any unused budget will be transferred to clusters with a smaller allocation weight.

\textbf{Text Budget Selection.}
We recommend running a minimum of $4 \times |\mathbb{F}|$ tests under the 3PA protocol.
In phase one, each fault needs one test to explore the potential interferences of the injection on the system.
In phase two, each fault is on average injected into two additional test cases, which enables the intra-group similarity calculation. We set a number of 2 instead of 1 in this phase just in case one fault is injected into a similar test workload as phase one. 
In phase three, we naturally allocate each fault one additional test. 
